\section{System and software architecture}
\label{sec:software-architecture}
The project is implemented as a batch-oriented Python application with separate data,
semantic, publication, and verification responsibilities. Figure~\ref{fig:system-architecture}
shows the complete system from external sources to user-facing artifacts. Figure~\ref{fig:software-architecture}
maps that flow to the repository's executable components. The diagrams are complementary:
the first explains the data life cycle, while the second explains the software boundaries.

\begin{figure}[htbp]\centering
\resizebox{\textwidth}{!}{\begin{tikzpicture}[
  node distance=7mm and 8mm,
  layer/.style={draw=accent!80!black, rounded corners=3pt, fill=accent!5, align=center, font=\small, text width=2.55cm, minimum height=13mm},
  source/.style={draw=orange!70!black, rounded corners=3pt, fill=orange!10, align=center, font=\small, text width=2.55cm, minimum height=13mm},
  artifact/.style={draw=blue!70!black, rounded corners=3pt, fill=blue!6, align=center, font=\small, text width=2.7cm, minimum height=13mm},
  arrow/.style={-{Stealth}, thick, draw=accent!80!black},
  sidearrow/.style={-{Stealth}, dashed, thick, draw=black!55}
]
\node[source] (sources) at (0,0) {Wikidata\\Vietnamese Wikipedia\\reference tables};
\node[source] (bronze) at (3.5,0) {Bronze\\cached responses\\source snapshots};
\node[layer] (silver) at (7,0) {Silver\\normalized records\\conflict log};
\node[layer] (transform) at (10.5,0) {RDF transformation\\persistent URI registry\\RDFLib};
\node[layer] (semantic) at (14,0) {Ontology + rules\\OWL 2 RL\\link graph};

\node[layer] (quality) at (10.5,-2.35) {Quality gates\\schema checks\\reasoning diagnostics};
\node[artifact] (release) at (14,-2.35) {Validated Gold release\\asserted RDF + selected inference\\SHACL and fingerprints};
\node[artifact] (static) at (10.5,-4.75) {Static Linked Data\\HTML, Turtle, JSON-LD\\canonical resource IRIs};
\node[artifact] (fuseki) at (14,-4.75) {Apache Jena Fuseki\\SPARQL + Graph Store\\TDB2 persistence};
\node[artifact] (app) at (17.5,-4.75) {Web application\\read-only query guard\\human-facing views};

\draw[arrow] (sources)--(bronze);
\draw[arrow] (bronze)--(silver);
\draw[arrow] (silver)--(transform);
\draw[arrow] (transform)--(semantic);
\draw[arrow] (semantic)--(release);
\draw[arrow] (quality)--(release);
\draw[sidearrow] (transform)--(quality);
\draw[arrow] (release.south west) to[bend right=12] (static.north east);
\draw[arrow] (release)--(fuseki);
\draw[arrow] (release.south east) to[bend left=12] (app.north west);

\node[font=\scriptsize, align=center] at (1.75,0.82) {acquire};
\node[font=\scriptsize, align=center] at (5.25,0.82) {integrate};
\node[font=\scriptsize, align=center] at (8.75,0.82) {map};
\node[font=\scriptsize, align=center] at (12.25,0.82) {enrich};
\end{tikzpicture}
}
\caption[Overall VN-Edu LOD architecture]{Overall architecture of VN-Edu LOD. Source material is
collected and normalized before RDF generation; ontology enrichment and quality gates produce a
validated release that can be published statically, queried through the application, or loaded
into Apache Jena Fuseki.}
\label{fig:system-architecture}
\end{figure}

\begin{figure}[htbp]\centering
\resizebox{\textwidth}{!}{\begin{tikzpicture}[
  node distance=7mm and 8mm,
  box/.style={draw=accent!80!black, rounded corners=3pt, fill=accent!5, align=center, font=\small, text width=3.1cm, minimum height=12mm},
  data/.style={draw=blue!70!black, rounded corners=3pt, fill=blue!7, align=center, font=\small, text width=3.1cm, minimum height=12mm},
  edge/.style={-{Stealth}, thick, draw=accent!80!black},
  aux/.style={-{Stealth}, dashed, thick, draw=black!55}
]
\node[box] (cli) at (0,0) {CLI orchestrator\\\texttt{run\_all.py}};
\node[box] (web) at (4.2,0) {Flask / Waitress\\Web UI and /sparql};
\node[box] (rdfdocs) at (8.4,0) {Static publisher\\HTML + RDF documents};
\node[box] (fuseki) at (12.6,0) {Fuseki adapter\\safe load and query};

\node[box] (collect) at (0,-2.2) {Collection and cache\\Wikidata/Wikipedia};
\node[box] (integrate) at (4.2,-2.2) {Integration\\identity and field policy};
\node[box] (transform) at (8.4,-2.2) {Transformation\\RDFLib + URI minter};
\node[box] (reason) at (12.6,-2.2) {Reasoning and validation\\owlrl + SHACL};

\node[data] (raw) at (0,-4.4) {Bronze snapshots};
\node[data] (silver) at (4.2,-4.4) {Silver JSON\\URI registry};
\node[data] (rdf) at (8.4,-4.4) {Gold RDF components};
\node[data] (reports) at (12.6,-4.4) {Release reports\\metrics and fingerprints};

\draw[edge] (cli)--(collect);
\draw[edge] (collect)--(integrate)--(transform)--(reason);
\draw[edge] (web)--(rdfdocs)--(fuseki);
\draw[edge] (collect)--(raw);
\draw[edge] (integrate)--(silver);
\draw[edge] (transform)--(rdf);
\draw[edge] (reason)--(reports);
\draw[aux] (rdf) to[bend left=18] (rdfdocs);
\draw[aux] (reports) to[bend right=18] (fuseki);

\node[font=\scriptsize, align=center] at (6.3,0.78) {application layer};
\node[font=\scriptsize, align=center] at (6.3,-1.48) {domain services};
\node[font=\scriptsize, align=center] at (6.3,-3.68) {versioned artifacts};
\end{tikzpicture}
}
\caption[Software architecture and repository components]{Software architecture. The CLI orchestrator
coordinates domain services, while versioned artifacts connect the services to the static publisher,
web application, and Fuseki adapter. Dashed arrows indicate artifact hand-off rather than direct
runtime calls.}
\label{fig:software-architecture}
\end{figure}

\subsection{Component responsibilities}
The collection layer obtains structured records and cached source responses. The integration
layer applies identity, field-precedence, normalization, and conflict policies. The RDF
transformer assigns or reuses resource IRIs, emits typed literals and object links, and keeps
the asserted graph separate from inferred output. The semantic layer applies the local ontology,
selected OWL/RDF rules, external links, and qualified observations. The release layer runs
raw-data checks, OWL diagnostics, SHACL, metric generation, graph-isomorphism checks, and
fingerprint validation before a serving artifact is accepted.

The web application is a consumer of the release rather than the source of truth. Its local
RDFLib backend is useful for an isolated read-only view; the optional Fuseki adapter provides
the tested HTTP/TDB2 path. The static publisher produces resource pages and RDF downloads for
Linked Data clients. This separation makes it possible to test the transformation and semantic
artifacts without requiring a running web server.

\subsection{Technology choices}
Python and RDFLib were retained for transparent transformation and easy regression testing.
OWL API and Protégé are used to inspect the ontology and check its OWL 2 RL profile. GraphDB
is useful for interactive repository exploration and visualization, but its layout is not the
formal specification. Apache Jena Fuseki 6.2.0 is used as the RDF HTTP server and TDB2
persistence target in the end-to-end acceptance workflow. Silk Framework 3.6.0 is isolated
as a candidate-linking experiment so that uncertain matches cannot silently enter the release.

\begin{table}[htbp]\centering\small
\caption{Main software components and their boundaries.}\label{tab:software-components}
\begin{tabularx}{\textwidth}{p{3.2cm}p{4.2cm}X}\toprule
Component & Responsibility & Boundary \\\midrule
\texttt{run\_all.py} and pipeline scripts & Orchestrate collection, integration, transformation, reasoning and publication & Batch execution; not a long-running RDF server \\
RDFLib and Python validators & Create RDF, apply local policies, inspect graphs, and run regression checks & Does not replace formal OWL profile checking or production load testing \\
Protégé/OWL API & Inspect ontology axioms and profile membership & Does not decide whether source facts are true \\
Flask/Waitress application & Read-only UI and guarded SPARQL access & Uses bounded local execution; external TLS/proxy policy remains operational \\
Apache Jena Fuseki/TDB2 & Persistent RDF store and SPARQL/Graph Store protocol & Requires deployment hardening, authentication, backups, and measured capacity \\
\bottomrule
\end{tabularx}
\end{table}
